\section{Description of the System}
\label{sec:system}

TrTok is implemented by a parallel execution of several configurable
pipeline steps.

\subsection{RoughTokenizer}

The RoughTokenizer partitions the stream of input characters into
small, discrete chunks of non-blank characters called \newterm{rough
tokens}. The partitioning can be made more granular by user-defined
rules which specify positions at which the desired tokenization might
differ from the whitespace-induced one.

A location in the text may be marked as a \maysplit{} meaning that the
characters in the text preceding and following it may be parts of
different tokens even though they are not separated by whitespace
(e.g. we might wish to put a \maysplit{} between \example{``was''} and
\example{``n't''} in \example{``wasn't''}).

A location within a span of white characters might be labeled as a
\mayjoin{} signalling that the characters preceding and following the
whitespace area might be parts of the same token, as in the case of
spaces entered in long numbers for readability (e.g.
\example{``12~345''}).

Finally, a location in the text may be marked as a \maybreaksentence{}
if the characters preceding it and the characters following it might
belong to different sentences.

\begin{figure}
  \begin{center}
    \includegraphics[width=\textwidth / 2]{img/decisionpoints.eps}
    \caption{In this example, $\blacktriangle$ stands for \maysplit{},
      $\blacktriangledown$ for \mayjoin{} and $\bullet$ for
      \maybreaksentence{}. This is how the rough tokenization might
      turn out given some reasonable settings for tokenizing English.}
    \label{fig:decision-points}
  \end{center}
\end{figure}

See Figure~\ref{fig:decision-points} for an example of how these
potential tokenization operations can look like in one sentence. A
rough token is defined as a maximal sequence uninterrupted by
whitespace or a possible tokenization operation (the symbols
underneath the sentence in Figure~\ref{fig:decision-points}). For
example, in the sentence from Figure~\ref{fig:decision-points},
``was'', ``n'', the apostrophe and ``t'' are all individual rough
tokens. Note that the presence of a \may{} event only signifies the
possibility of a tokenization operation (splitting or joining of
tokens or sentences). Whether a token split, token join or sentence
break will occur is up to the Classifier.

The locations of these possible tokenization operations are determined
by user-defined rules which each consist of a pair of regular
expressions. The respective tokenization operation is signalled if the
characters leading up to and following a position match the regular
expressions in one of these rules.

If we look back at Figure~\ref{fig:decision-points}, we might imagine
more robust settings also placing a \maybreaksentence{} after the
apostrophe/single quote, while others might be more daring and not
place a \maybreaksentence{} after the point in ``2.0'', because it is
followed by a non-blank character.

TrTok collects these rules and generates a Quex program implementing a
fast FSM, which gets compiled and loaded in
(Quex\footnote{http://quex.sourceforge.net} is a fast and more
Unicode-friendly variation on the classic tools lex and flex for C++).

\subsection{FeatureExtractor}

The stream of rough tokens interleaved with potential tokenization
operations output by the RoughTokenizer is processed using the
FeatureExtractor. The FeatureExtractor annotates each rough token with
a bit vector signifying which of the user-defined feature predicates
hold for the rough token in question.

The features can be defined in two ways: either using a regular
expression or a list of rough tokens. For a feature defined using a
regular expression, a rough token is said to have that feature if and
only if the regular expression matches the entire rough token. In the
case of a feature defined using a list of rough tokens, a rough token
is said to have the feature if and only if it is on the list.

This way it is easy to specify features which try to analyze the shape
of rough tokens using regular expressions or to simply give a list of
all interesting tokens (e.g. words of a certain part of speech or
exceptions such as abbreviations).

\subsection{Classifier}

The Classifier is the other ``hard worker'' of the system (besides the
RoughTokenizer). Its job is to disambiguate the potential tokenization
operations identified by the RoughTokenizer, i.e. it decides whether a
\maysplit{} truly splits a word into two tokens, whether a \mayjoin{}
joins two words into one token and whether a \maybreaksentence{} truly
ends a sentence. It does so by consulting a maximum entropy classifier
for every location containing these potential tokenization operations.

The features passed to the classifier consist of the features of rough
tokens in the context surrounding the potential tokenization operation
and the presence of whitespace and potential tokenization operations
in the context area. The user is free to select the size of the
context area and which features from which rough tokens in the context
area should be passed to the classifier.

Features can also be clustered together into conjunct features which
provide a value for every combination of the constituent features'
values (this lets the trainer estimate a different parameter for
different combinations of the features' values, which is useful to
model the joint influence of some features).

The classifier then marks each location with a potential tokenization
operation as either a sentence boundary, token boundary or no boundary
(the location is inside a token). Using this classification, any
potential tokenization operations are finally disambiguated.

\subsection{OutputFormatter}

The OutputFormatter is the point at which the stream of rough tokens
is turned back into a character stream. This means that all the rough
tokens are concatenated and whitespace is inserted between them
depending on whether there originally was any whitespace between them
and on the tokenization operations which are to be carried out in the
space between them. Individual tokens end up being separated by a
single space character and sentences are separated by line breaks.
